\chapter{Conclusion générale}
\label{ch:conclusion}

\section{Bilan}

Le projet avait pour objet la réalisation d'une plateforme web de gestion pour le
laboratoire LIAS, remplissant simultanément la fonction de site institutionnel public et
celle de système d'information interne. Les 21 sections du cahier des charges sont
couvertes par des fonctionnalités opérationnelles, connectées à une base de données réelle
et vérifiées par un parcours complet de l'application.

Le résultat se mesure sur trois plans. Sur le plan fonctionnel, dix-huit modules couvrent
le cycle de vie des membres, la gouvernance, l'activité scientifique, les ressources
matérielles et documentaires, la communication interne et le pilotage annuel. Sur le plan
technique, le système compte 31 entités persistantes, 22 contrôleurs REST et 21 pages
d'interface, soit environ 23~000 lignes de code applicatif. Sur le plan de la validation,
55 tests automatisés s'exécutent sans échec et un audit fonctionnel a porté sur l'ensemble
des routes et des rôles.

La difficulté centrale du sujet --- la séparation stricte entre le \srt{} d'appartenance,
le \man{} de gouvernance et le \rol{} applicatif --- a été traitée par une décision de
conception prise tôt et tenue jusqu'au bout : le rôle n'est pas une donnée saisie mais une
donnée dérivée. Cette décision a eu des conséquences en cascade sur le modèle de données,
sur le service de gouvernance, sur la chaîne de sécurité et jusque sur la présentation des
fiches membres. Elle a également révélé sa propre limite --- l'existence de comptes
antérieurs non réconciliés --- que la procédure de réconciliation au démarrage est venue
combler.

Le bilan comporte aussi ce qui n'a pas été fait : le thème sombre, la mesure de couverture,
la couverture de test de trois modules, la campagne de performance et l'assainissement de
l'historique de gestion de version. Ces points sont énoncés au chapitre précédent et n'ont
pas été dissimulés dans la conclusion.

\section{Compétences acquises}

\paragraph{Sur le plan technique.} La conduite du projet a permis de pratiquer l'ensemble
de la chaîne Jakarta~EE telle qu'elle est mise en œuvre par Spring~Boot : mapping objet-
relationnel avec JPA, validation déclarative par Bean~Validation, traitement des requêtes
au niveau de l'API~Servlet, démarcation transactionnelle, et sécurité par annotation. Le
versionnement du schéma par migrations et l'exécution de l'ORM en mode strict de validation
ont imposé une rigueur sur le modèle de données qu'une génération automatique du schéma
aurait masquée.

\paragraph{Sur le plan de la conception.} La leçon la plus durable concerne la distinction
entre donnée saisie et donnée dérivée. Faire du rôle un champ modifiable aurait été plus
simple à écrire et aurait fonctionné --- jusqu'à la première incohérence entre le rôle
stocké et la réalité des mandats, incohérence impossible à détecter automatiquement. Le
choix inverse coûte un service supplémentaire et impose une procédure de réconciliation,
mais il rend une classe entière d'anomalies structurellement impossible.

\paragraph{Sur le plan de la validation.} Le projet a montré concrètement les limites de
chaque niveau de vérification. La compilation ne dit rien du comportement. Un code de
statut HTTP correct ne garantit pas une réponse correcte --- le défaut de sérialisation des
réponses d'erreur en est l'illustration la plus nette, puisqu'il a survécu à des dizaines
de vérifications portant sur le seul code retour. Seul le parcours réel de l'application,
prolongé par des tests d'intégration exerçant la chaîne complète, a permis de le mettre au
jour.

\paragraph{Sur le plan méthodologique.} La gestion des secrets constitue un enseignement à
part. Un identifiant de base de données placé dans un fichier de configuration au tout
début du projet y est resté jusqu'à la phase finale, et son retrait a exigé, en plus d'une
modification de code, une rotation effective du mot de passe. Une configuration de
développement est un fichier de code comme un autre.

\section{Apports du projet}

Pour le laboratoire, la plateforme apporte une mémoire institutionnelle interrogeable : la
succession des directions, les périodes d'affiliation, les décisions d'adhésion, les
documents et leurs versions successives cessent d'être dispersés dans des supports qui ne
conservent ni l'ordre ni l'auteur des décisions. La production automatisée du rapport
d'activité annuel, qui conditionne le budget, transforme un exercice de compilation manuel
en une agrégation reproductible.

Pour la formation, le projet a servi de terrain d'application à des notions habituellement
abordées séparément : persistance, sécurité, architecture en couches, modélisation UML et
validation. Leur mise en œuvre conjointe sur un domaine métier non trivial fait apparaître
des interactions qu'un exercice isolé ne révèle pas --- notamment entre la suppression
logique et l'intégrité référentielle, ou entre l'authentification sans état et la
révocation immédiate des droits.

\section{Perspectives}

À court terme, les limites recensées au chapitre~\ref{ch:resultats} constituent une feuille
de route immédiate : thème sombre, couverture de test objectivée, extension des tests aux
modules non couverts, rafraîchissement transparent du jeton, purge de l'historique de
gestion de version.

À plus longue échéance, la perspective la plus cohérente avec la spécialité du Master
consiste à exploiter analytiquement les données que la plateforme accumule. Le principe de
conservation intégrale, adopté ici pour satisfaire une exigence de traçabilité
administrative, a un effet de bord favorable : chaque table d'historique constitue une
série temporelle. Les effectifs par statut, la durée des mandats, la productivité
scientifique par équipe ou la saisonnalité des activités sont déjà présents dans le
schéma ; il ne manque qu'une couche de restitution pour les rendre exploitables.

Le système livré n'est pas un prototype : il fonctionne sur une base réelle, applique une
politique de sécurité vérifiée et conserve l'intégralité de l'historique qu'il produit. Il
constitue, à ce titre, une base sur laquelle il est possible de construire.
